\section{Considerações finais e próximos passos}
\label{sec:conclusao}

O protótipo 2D (\code{IBHMInterfaceElementBuilder}) cumpre o que a apresentação propôs: com quatro operadores locais,
escolhidos apenas pelo estado do vértice, insere elementos de interface em qualquer ordem, sobre malhas híbridas que já
contêm elementos de interface, em uma estrutura compacta baseada em índices e com memória reservada antecipadamente e de
forma exata. A verificação mostrou, em malhas de triângulos, quadriláteros, polígonos e com furos, que o resultado é a
topologia esperada (uma cópia por componente da estrela cortada), válido e independente da ordem e do modo de
inserção, com contagem exata de vértices. Desde a v0.3 foram corrigidos todos os defeitos que afetavam o resultado:
cópias com as coordenadas do original, entrada validada, contagem exata, vértice de borda com várias arestas, laço
infinito latente e sentinelas protegidos por asserções. Os princípios do algoritmo têm testes automatizados próprios, e
o teste de mutação mostrou que eles pegam erros que deixam o resultado válido mas errado. O que resta em aberto é relato
à aplicação (I7), desempenho (I9) e memória temporária na construção (I12).

\subsection{Um invariante que resume o algoritmo}

Uma forma compacta de enunciar o que os operadores constroem, e que pode orientar tanto os testes quanto a extensão 3D,
é a seguinte:

\begin{nota}[Invariante]
Depois de inseridas as arestas de um conjunto $S$, cada vértice original $v$ tem tantas cópias quantos são os
componentes conexos da sua estrela depois de removidas as arestas de $S$. Cada componente usa uma única cópia.
\end{nota}

Para um vértice interno com $k$ arestas selecionadas, isso dá $1$ cópia se $k\le1$ (a ponta de trinca, que fica
colapsada) e $k$ cópias se $k\ge2$; para um vértice de borda, $k+1$ cópias. O buraco que aparece com $k\ge3$ é a
consequência geométrica de ligar $k$ cópias apenas por elementos de espessura nula. Como a partição em componentes não
depende da ordem, qualquer algoritmo correto chega ao mesmo resultado, e é isso que explica a independência da ordem
observada. Os operadores são uma forma \emph{incremental} e local de manter esse invariante.

\subsection{Rumo ao 3D}
\label{sec:3d}

Algumas observações, como hipóteses de trabalho para a próxima etapa:
\begin{itemize}
\item Em 3D a entrada são faces, e o elemento de interface é um prisma (faces triangulares) ou hexaedro (faces
quadrilaterais) de espessura nula. A duplicação deixa de acontecer só nos vértices: as \emph{arestas} das faces
selecionadas também precisam ser duplicadas. Uma aresta compartilhada por $k$ faces selecionadas se comporta como um
vértice 2D com $k$ arestas selecionadas, e girar em torno de uma aresta em 3D (de célula em célula pelas half-faces
opostas) é o análogo exato de girar em torno de um vértice em 2D. Os operadores 2D devem reaparecer, quase sem
mudanças, como operadores \emph{de aresta}.
\item A estrela de um vértice em 3D não é mais um ciclo, e sim uma superfície fechada (ou com borda) de células. Dividir
essa estrela exige uma busca em largura pelas células que não atravessa faces selecionadas, e não uma rotação. O
invariante acima continua valendo e dá o critério: uma cópia por componente.
\item O estado colapsado tem análogos em 3D (faces laterais degeneradas do prisma), e o buraco das junções vira
``tubos'' ao longo das arestas de junção.
\item O sentinela \COLLAPSED, a pré-contagem e a estratégia contar/reservar/validar se transportam diretamente para a
3D-IBHM.
\end{itemize}

\subsection{Próximos passos sugeridos}

\begin{enumerate}
\item \textbf{Extensão 3D} com os testes como base. O oráculo da topologia esperada (\cref{sec:testes}) se transporta
quase sem mudanças: uma cópia por componente da estrela de cada vértice (e de cada aresta) cortada pelas faces
selecionadas. O gerador de malhas aleatórias, as verificações de estrutura (gabarito do elemento, decisão pelo estado,
contrato de cada operador, localidade) e a equivalência entre o operador particular e o geral (Split/Expand) podem ser
reescritos para 3D antes dos operadores, servindo de especificação executável.
\item \textbf{Relatar resultados} (I7): uma entrada por aresta de entrada, com o elemento criado ou o motivo da
rejeição; a biblioteca deixa de imprimir. Retomar junto o B10.
\item \textbf{Transferência de estado}: expor o mapa ``vértice novo $\to$ vértice de origem'', no mesmo formato. Numa
análise adaptativa, os campos nodais (deslocamentos, poropressões etc.) dos vértices novos precisam ser copiados do
vértice de origem; o construtor conhece essa informação no momento da criação (as coordenadas já são copiadas, I4).
\item \textbf{Simplificar os operadores} em torno de uma primitiva comum de ``divisão de leque'' (\cref{sec:relacoes}),
mantendo o Split Element como otimização.
\item \textbf{Desempenho e memória} (I9, I12): manter o vetor $VH[\,]$, reaproveitar o intervalo do elemento corrente
nas rotações e montar as opostas sem um vetor por vértice.
\item \textbf{Integração com o GeMA}.
\end{enumerate}
